\section{What the experiments establish}
\label{sec:experiments}
The work formulas tell us what each method must do, but we need measurements to find out
whether the scoring it avoids pays for the search it adds. We compare real embeddings first,
then use constructed inputs to isolate how much the query's strong coordinates matter.
The saved source record identifies the runs behind these comparisons~\cite{local}.

\subsection{Make the comparison fair}
A, B and C receive the same documents and floating-point queries, use the same native
weighted binary scorer, and return the top 100 with score ties broken by document ID.
When we measure recall here, we measure recovery of that exact binary ranking; these
comparisons use neither a float reranker nor human relevance judgments.

The machine is an Apple M5 Max with 128 GiB of physical memory, while each worker runs
under a separate 32 GB decimal memory budget using one CPU thread and one query at a time.
A fresh process builds a configuration, warms it, and then measures routing plus search.
The outer timer includes the call from Python into C++, but cached-vector measurements
exclude text encoding. Repeating whole process blocks helps expose timing variation;
the recorded power settings do not prove that CPU frequency remained fixed.

We tune parameters on one set of query IDs, freeze the choices, and evaluate them on different
IDs, allowing each method its own fastest tested setting that meets the recall requirement.
For A, the tested options include intermediate probe counts and direct routing by sign keys.
For B, they include large leaves, so it can score a selected cluster without unnecessary
splits. If we left these options out, a poor baseline or a forced traversal could decide
the result before the timing begins.

The real study retains two ways to select settings: one accepts a setting when its average
tuning recall reaches the target, while the conservative rule requires a lower estimate from
repeated resampling of tuning queries to reach it. The conservative rule leaves more margin,
but cannot promise recall on future queries. To establish a speedup, both methods must also
meet the target and the paired 95\% interval for their time ratio must lie above one. Queries
and whole process blocks are resampled together; timing repeats do not count as additional
independent queries.

\subsection{Real embeddings did not establish an alternative speedup}
For the real study, the documents are 256-bit codes from Nomic embeddings of MS MARCO passages~\cite{nomic,msmarco}.
After combining 6,821 tuning settings, the frozen choices were evaluated on 200 query IDs
excluded from the earlier 300-query collection. On this sample, all 36 conservative target
points met their requested recall, compared with only 7 of 36 points chosen by average
tuning recall.

In the conservative results at one million documents below, each row sets a common minimum
recall requirement and shows each method's median routing-plus-search time. If we compare
methods within a row, their achieved recall also happens to be equal.

\begin{center}
\begin{tabular}{@{}rrrrr@{}}
\toprule
Target & A (ms) & B (ms) & C (ms) & Recall (\%)\\
\midrule
80\% & 0.675 & 0.699 & 0.751 & 82.53\\
90\% & 1.634 & 1.686 & 1.880 & 91.73\\
95\% & 3.360 & 3.430 & 3.706 & 95.47\\
99\% & 9.580 & 9.965 & 11.013 & 99.25\\
\bottomrule
\end{tabular}
\end{center}

\fig{real-tradeoff}{Four frozen recall-target choices on one million real passage embeddings and 200 independent query IDs. Lines connect the tested settings for readability; they do not predict intermediate settings. At each shown recall, A has the smallest median.}{.9}

Across the real study, none of the 48 evaluated comparisons established a B or C speedup
under the stated rule. Although fewer documents received full scores, that saving did not
produce a supported reduction in complete retrieval time. This conclusion applies to this
data, score, implementation and tested parameter range.

A later control gives C a one-bit key, A's routing layout, and unlimited candidate and lookup
budgets, while A and B keep their frozen settings. This control uses 100 further query IDs,
excluding all 500 previously used IDs. A, B and C measured 9.541, 9.855 and 10.545 ms,
respectively, and all three achieved 98.95\% recall, below the requested 99\%. That remains
a miss; we do not round it into a pass. The sample is separate from the earlier successful
200-query evaluation, and no parameter was changed in response to its outcome.

\subsection{Constructed data shows when skipping scores can help}
With constructed inputs, we can isolate the effect of keeping the strong coordinates fixed
or changing them with each query. Document signs are independent and balanced across 256
coordinates. Each query has 13 strong coordinates of magnitude one, with magnitude 0.0001
on the remaining coordinates. A document matching every strong sign therefore beats any
that mismatches one, regardless of the weak signs. The fixed case keeps the strong positions
the same; in the changing case, each query chooses different strong positions uniformly
without replacement. These vectors are constructed, not measured text embeddings.

The final follow-up uses 40 tuning queries and 20 separate evaluation queries, with three
process blocks and three timing repeats. It adds neighboring key widths, cluster counts and
leaf sizes, and then remeasures earlier choices. At one million documents and the 99\%
requirement, the results below show that all methods reach 100\% recall except A in the
changing case, at 99.95\%.

\begin{center}
\begin{tabular}{@{}lrrr@{}}
\toprule
Strong positions & A (ms) & B (ms) & C (ms)\\
\midrule
Fixed & 0.0140 & 0.0160 & 0.0128\\
Changing & 24.9235 & 0.1446 & 14.6301\\
\bottomrule
\end{tabular}
\end{center}

With fixed positions, A selects two of 16,384 clusters by direct routing, while C uses one
physical cluster and a 13-bit key. C's paired speedup estimate is 1.098, with a 95\% interval
of 1.066--1.147, and the time difference is about 1.25 microseconds. Both methods select
documents through the same strong signs, but C combines key enumeration and scoring in one
native call, whereas A routes and scans through separate calls. The small measured advantage
therefore concerns this implementation of closely related candidate selection.

If the strong positions change, a fixed key cannot consistently select them. B can follow
them because it orders coordinates using each query's magnitudes; its selected setting uses
one cluster, unlimited traversal and a leaf size of 160. A opens 255 of 256 centroid clusters,
and B's paired speedup estimate is 172.3, with a 95\% interval of 169.1--177.2. C uses a
14-bit key and has a lower median than A, but its paired interval against A includes one,
so that median alone does not establish an advantage.

We can see the reason for the leaf choice in the occupancy calculation: at one million rows,
the ideal 13-bit group has mean size about 122 and standard deviation about 11. A threshold
of 128 can force another split on a weak coordinate, whereas a threshold of 160 leaves more
room to score that group. The experiment checks this suggestion from the distribution;
it does not establish that 160 is best for arbitrary embeddings, sizes or machines.

\subsection{Check the answer, the work, then the predicted time}
Before checking work or timing predictions, we need an independently calculated answer.
The checks calculate scores directly from document signs and sort by score and ID; larger
runs count strong and weak mismatches from packed bits to avoid a huge expanded matrix.
Every complete native call in the scaling check returned the independent reference ranking.

To check the predicted work, we compare it with native counters. Across the earlier
controlled tuning and evaluation runs, 447 applicable per-query predictions matched exactly,
including full-width split and leaf visits, scored rows, key attempts including empty
lookups, and tracked mask peaks where applicable. Queries outside the formulas' stated
conditions remain in the measurements, but do not count as successful predictions.

Even with correct counts, we cannot assign every row the same time cost. In a
16-million-document buffer, scoring about 1,958 scattered rows took 0.159 ms with reused
candidates and 0.455 ms with changing candidates. Candidate generation stayed outside the
timer, but still changed which memory had been accessed recently. Fits trained on one- and
four-million-row buffers underestimated a 16-million-row check by as much as 45.2\% for
reused candidates and 27.6\% for changing candidates, and both failed fits remain saved.
The largest row store was about 640 MB, so these checks do not validate access costs for
a 40 GB store.

The later complete-call model uses one latency median per corpus size, correcting an earlier
fit that gave fast individual queries too much weight when predicting a median. The revised
model also accounts for result-selection work in C's small postings. For B, the leaf threshold
grows as $L=160N/10^6$, so a larger corpus keeps the intended strong-coordinate depth
instead of forcing more cuts.

Because the revised four-million-row check had already been inspected, it serves as a
diagnostic check. The coefficients were frozen before the separate eight-million-row run
below, where a positive error means the prediction was slower than the measurement.

\begin{center}
\begin{tabular}{@{}llrrr@{}}
\toprule
Positions & Method & Measured (ms) & Predicted (ms) & Error\\
\midrule
Fixed & A & 197.029 & 203.016 & $+3.0\%$\\
Fixed & B & 0.854 & 0.884 & $+3.6\%$\\
Fixed & C & 0.0481 & 0.0518 & $+7.7\%$\\
Changing & A & 196.664 & 197.816 & $+0.6\%$\\
Changing & B & 0.900 & 0.859 & $-4.5\%$\\
Changing & C & 202.767 & 207.086 & $+2.1\%$\\
\bottomrule
\end{tabular}
\end{center}

All six checks met the earlier 20\% criterion, testing transfer to a larger corpus under
the same prescribed workload and existing query IDs. These model-check configurations are
separate from the tuned comparisons in the preceding table. Although the largest
observed error was 7.7\%, that does not give us a guaranteed bound for future queries or
a billion documents.

\subsection{Measure actual text requests}
When a user supplies text, retrieval covers only part of the wait. The final check therefore
runs 20 existing independent-query texts through the cached local Nomic model on MPS,
normalizes each vector, routes it and calls native search, with a batch size of one.
Three blocks and two repeats give 120 requests per method, or 360 requests in total. Each
sequence shape is warmed, and the request timer excludes model loading, index building
and reference creation.

Using the earlier frozen settings, all three methods achieve 99.45\% binary recall on this
sample. The table gives the stage medians alongside the separately measured request total.

\begin{center}
\begin{tabular}{@{}lrrrr@{}}
\toprule
 & Encode (ms) & Retrieve (ms) & Total (ms) & Total p95 (ms)\\
\midrule
A & 6.488 & 9.406 & 16.313 & 19.201\\
B & 6.500 & 9.651 & 16.303 & 19.217\\
C & 7.543 & 11.201 & 18.912 & 21.898\\
\bottomrule
\end{tabular}
\end{center}

The median is the middle time, and p95 is the 95th-percentile time. The tiny difference
between the A/B totals does not establish a speedup. If we add the stage medians, we need
not get the total median, because the request with the middle encoding time need not have
the middle retrieval time. We must measure each request's total before taking its median.

These timings apply to real Nomic text queries. Since the constructed strong-coordinate
vectors do not come from this encoder, adding its time to their search measurements would
not give us a measured text-to-results pipeline.

% Source map (paths relative to repository root; final boundary results supersede 07c timings):
% Fairness and real study: reports/search-math/sections/07b-implemented-study.tex;
% results/independent-evaluation-2026-09-10/report/{targets,comparisons}.csv.
% Real numeric table: reports/search-math/evidence/independent-results-table.tex.
% Final selected evidence: reports/search-math/evidence/final-study.json; all source hashes checked.
% Later real control: results/one-bit-evaluation-2026-09-11/report/targets.csv.
% Constructed distribution: results/controlled-boundary-adaptive-2026-09-11/config.toml.
% Final controlled timings/settings: results/controlled-boundary-{fixed,adaptive}-2026-09-11/
% evaluation/report/{targets,comparisons}.csv. Fixed C interval checked directly in comparisons.csv.
% Conditional counts: results/controlled-{fixed,adaptive}-2026-09-10/{tuning,evaluation}/work-checks.json;
% checked totals 193+120+92+42=447; zero saved mismatches.
% Failed fits: results/{buffer-analysis,buffer-fresh-analysis}-2026-09-11/leaf-model.json.
% Model check: results/native-scale-confirmation-2026-09-11/frozen-model-check.json;
% explanation: reports/search-math/sections/07d-cost-validation.tex.
% Text scope and values: results/text-queries-2026-09-11/{configuration,summary}.json.
